\section{Постановка задачі}
\label{sec:problem}

\subsection{Геометрія та вихідна конфігурація}
\label{subsec:geometry}

В умовах плоскої деформації розглядається кусково-однорідне тіло з
ламаною межею поділу двох різних пружних ізотропних матеріалів
(рис.~\ref{fig:failure-stages}). Матеріал
$i$ характеризується модулем Юнга $E_i$ і коефіцієнтом Пуассона $\nu_i$
($i=1,2$). Кут зламу межі поділу, відлічений у матеріалі~2, дорівнює
$2\alpha$, де $\alpha$ --- півкут зламу,
\begin{equation*}
  0<\alpha<\pi, \qquad \alpha\ne\frac{\pi}{2}.
\end{equation*}
Значення $\alpha=\pi/2$ відповідає плоскій межі поділу і надалі
розглядається лише як граничний випадок.

Початок полярної системи координат $(r,\theta)$ розміщено в кутовій
точці $O$. Завдяки симетрії відносно бісектриси кута достатньо розглядати
півплощину $0\leq\theta\leq\pi$: матеріал~2 займає сектор
$0<\theta<\alpha$, матеріал~1 --- сектор $\alpha<\theta<\pi$, а межі
поділу відповідає промінь $\theta=\alpha$. Компоненти переміщення
позначаємо через $u_r,u_\theta$, а компоненти напруження --- через
$\sigma_\theta,\tau_{r\theta}$.

Тіло навантажене так, що локальне поле біля точки $O$ є симетричним
відносно бісектриси. На межі поділу його сингулярна частина має вигляд
\begin{equation*}
\begin{aligned}
  \tau_{r\theta}(r,\alpha)
    &= Cg_1 r^{\lambda_0}+o\!\left(r^{\lambda_0}\right),\\
  \sigma_\theta(r,\alpha)
    &= Cg_2 r^{\lambda_0}+o\!\left(r^{\lambda_0}\right),
    \qquad r\to0.
\end{aligned}
\end{equation*}
Тут $C$ --- амплітуда локального сингулярного поля, яка враховує
вплив зовнішнього навантаження і глобальної геометрії тіла;
$g_1,g_2$ --- функції кута $\alpha$ та пружних характеристик матеріалів.
Показник $\lambda_0\in(-1,0)$ визначає сингулярність інтактної кутової
точки. Під $\lambda_0$ розуміємо найменший дійсний корінь рівняння
$D_0(-1-\lambda_0)=0$ у відкритому інтервалі $(-1,0)$, неперервно
пов'язаний із фізичною гілкою розв'язку; функції $D_0,g_1,g_2$ наведено в
додатку~\ref{app:asymptotics}.

У цій роботі вихідною конфігурацією є тіло з уже наявною заданою
стаціонарною V-подібною міжфазною тріщиною зсуву, утвореною двома
симетричними гілками, що виходять із кутової точки $O$ уздовж ламаної
межі поділу. Береги тріщини контактують без тертя. Довжина кожної
гілки дорівнює $l$, тому сумарна довжина двох гілок становить $2l$;
надалі $l$ називаємо півдовжиною тріщини (довжиною її гілки).
Контактний стан узгоджується з умовою
\begin{equation*}
  Cg_2<0,
\end{equation*}
за якої нормальне напруження на межі поділу є стискальним. У цій
постановці $l$ є заданим і фіксованим параметром. Рівень зовнішнього
навантаження вважаємо недостатнім для втрати локальної стійкості біля
зовнішніх вершин тріщини. Міжфазне поширення з цих вершин не входить до
локальної задачі біля $O$ і нижче розглядається лише як конкурентний
механізм руйнування.

Кутова точка заданої міжфазної тріщини залишається концентратором
напружень зі степеневою особливістю~\cite{NazarenkoKipnis2022Stress}.
Це поле може підтримувати в одному з матеріалів малу зону
передруйнування, подальша трансформація якої у тріщину нормального
відриву описує ініціювання відгалуження біля кутової точки. За прийнятої симетрії
геометрії та навантаження природним напрямком такої зони є бісектриса кутового сектора,
зайнятого цим матеріалом. Для бісектрисного напрямку локальна складова
другої моди в кінці зони зникає за симетрією, що узгоджується з принципом
локальної симетрії~\cite{BarenblattCherepanov1961,GoldsteinSalganik1974}. У цій роботі зони з довільною орієнтацією не
розглядаються, тому бісектрисний напрямок приймається як частина
симетричної моделі, а не визначається шляхом окремої оптимізації напряму.

Прийнята нумерація напрямків є
\begin{equation*}
  \beta_1=\pi,\qquad \beta_2=0.
\end{equation*}
Індекс $i$ у величинах $\beta_i,d_i,\sigma_i,Q_i$ надалі завжди означає
номер матеріалу.

Зону передруйнування в матеріалі~$i$ моделюємо лінією розриву нормального
переміщення довжиною $d_i$, на якій діє стале відривне напруження
$\sigma_i>0$. На рис.~\ref{fig:failure-stages} суцільною лінією показано зону
в матеріалі~1, а штриховою --- альтернативну зону в матеріалі~2. Для
кожного допустимого напряму розрахунок визначає сім'ю квазістатичних
рівноважних станів зони; її розкриття та енергетична характеристика
надалі використовуються для формулювання критеріїв ініціювання
відгалуженої тріщини.

\begin{figure}[htbp]
  \centering
  % Vector redraw of the author's original Figure 1.
% The figure is intentionally kept as TikZ so that all labels are typeset
% by LaTeX and remain vector objects in the compiled manuscript PDF.
\begin{tikzpicture}[x=1cm,y=1cm,>=Stealth,line cap=round,line join=round]
\tikzset{
  panel/.style={draw=black,line width=.6pt,rounded corners=6pt},
  interface/.style={draw=black,line width=.7pt},
  axis/.style={draw=black,line width=.45pt},
  crack/.style={draw=black,double=white,double distance=2.0pt,line width=.45pt},
  processzone/.style={draw=black,line width=2.0pt},
  alternativezone/.style={draw=black,line width=1.2pt,dashed}
}

% Common panel geometry.  The interface rays are symmetric about the
% horizontal bisector of material 2, as in the author's original sketch.
\newcommand{\failurepanel}[2]{%
  \begin{scope}[xshift=#1cm]
    \draw[panel] (-2.20,-1.55) rectangle (2.20,1.55);
    \coordinate (O) at (0,0);
    \coordinate (Iu) at ($(O)+(118:2.05)$);
    \coordinate (Il) at ($(O)+(242:2.05)$);
    \draw[interface] (O)--(Iu);
    \draw[interface] (O)--(Il);
    \draw[axis] (-2.05,0)--(2.05,0);

    % Polar coordinates and the half-angle alpha.
    \draw[->,line width=.65pt] (O)--++(27:1.62)
      node[above right=-1pt,font=\small] {$r$};
    \draw[->,line width=.5pt] (0.58,0)
      arc[start angle=0,end angle=27,radius=.58];
    \node[font=\small,fill=white,inner sep=.7pt] at (0.93,0.18) {$\theta$};
    \draw[->,line width=.5pt] (0.83,0)
      arc[start angle=0,end angle=118,radius=.83];
    \node[font=\small,fill=white,inner sep=.7pt] at (0.16,0.88) {$\alpha$};

    \node[font=\small,fill=white,inner sep=.7pt] at (0.18,-0.20) {$O$};
    \node[font=\small] at (-1.45,-0.92) {$E_1,\,\nu_1$};
    \node[font=\small] at ( 1.30,-0.92) {$E_2,\,\nu_2$};
    #2
  \end{scope}%
}

% (a) Intact broken interface.
\failurepanel{0}{%
  \node[font=\small] at (0,-1.88) {(a)};
}

% (b) Symmetric interfacial shear cracks of length l.
\failurepanel{5.00}{%
  \draw[crack] (O)--++(118:1.30);
  \draw[crack] (O)--++(242:1.30);
  \node[font=\small] at (-0.77, 1.18) {$l$};
  \node[font=\small] at (-0.77,-1.18) {$l$};
  \node[font=\small] at (0,-1.88) {(b)};
}

% (c) Mode-I process-zone alternatives along the two material bisectors.
\failurepanel{10.00}{%
  \draw[crack] (O)--++(118:1.30);
  \draw[crack] (O)--++(242:1.30);
  \node[font=\small] at (-0.77, 1.18) {$l$};
  \node[font=\small] at (-0.77,-1.18) {$l$};
  \draw[processzone] (O)--(-1.35,0);
  \draw[alternativezone] (O)--(1.35,0);
  \node[font=\small,below=2pt] at (-0.75,0) {$d_1$};
  \node[font=\small,below=2pt] at ( 0.75,0) {$d_2$};
  \node[font=\small] at (0,-1.88) {(c)};
}
\end{tikzpicture}

  \caption{Локальні конфігурації біля кутової точки:
  (a) інтактна ламана межа поділу; (b) задана стаціонарна V-подібна
  міжфазна тріщина зсуву з двома гілками довжиною $l$ кожна
  (сумарна довжина гілок $2l$); (c) зона передруйнування довжиною $d_1$ у матеріалі~1 та
  альтернативна зона довжиною $d_2$ у матеріалі~2 (штрихова лінія).
  Панелі зіставляють відповідні локальні конфігурації і не задають
  послідовності їх реалізації.}
  \label{fig:failure-stages}
\end{figure}

\subsection{Локальна задача та крайові умови}
\label{subsec:boundary-conditions}

Нехай $L$ --- характерний зовнішній розмір тіла. Приймаємо розділення
масштабів
\begin{equation*}
  d_i\ll l\ll L
\end{equation*}
і нехтуємо тертям між берегами міжфазної тріщини. У межах цієї
квазістатичної постановки інтенсивність зовнішнього поля $C$ може
змінюватися, тоді як півдовжина тріщини $l$, а отже й сумарна довжина
її гілок $2l$, залишаються фіксованими. Зміна навантаження змінює
рівноважну довжину зони $d_i$ та пов'язані з нею локальні
характеристики. У локальній задачі
кусково-однорідне тіло замінюємо нескінченною площиною, міжфазні тріщини
вважаємо півнескінченними, а вплив їхньої скінченної довжини та зовнішнього
навантаження задаємо умовою зрощування асимптотичних розв'язків~\cite{Nayfeh1973Perturbation}.
Така постановка є однобічним локальним наближенням: відоме поле
стаціонарної міжфазної тріщини задає навантаження для внутрішньої задачі
біля $O$, тоді як зворотний вплив малої зони передруйнування на поля
біля зовнішніх вершин тріщини не визначається.

Для величини $f$, заданої по обидва боки променя $\theta=\alpha$,
використовуємо позначення
\begin{equation*}
  \jump{f}
  =f(r,\alpha+0)-f(r,\alpha-0).
\end{equation*}
Умови симетрії, рівноваги і контакту на міжфазній тріщині мають вигляд
\begin{equation}
\label{eq:interface-boundary-conditions}
\begin{aligned}
  &\theta\in\{0,\pi\}: &&\tau_{r\theta}=0,\\
  &\theta=\alpha: &&
    \jump{\sigma_\theta}=0,
    \quad \jump{\tau_{r\theta}}=0,
    \quad \tau_{r\theta}=0,
    \quad \jump{u_\theta}=0.
\end{aligned}
\end{equation}
Рівність $\jump{u_\theta}=0$ описує контакт берегів у
нормальному напрямку, тоді як $\tau_{r\theta}=0$ відповідає відсутності
тертя.

На бісектрисі кутового сектора матеріалу, в якому розташована зона
передруйнування, задаємо змішані умови. Для позначення іншого матеріалу
використовуємо індекс $j=3-i$:
\begin{equation}
\label{eq:process-zone-boundary-conditions}
\begin{aligned}
  &\theta=\beta_i,\quad 0<r<d_i:
    &&\sigma_\theta(r,\beta_i)=\sigma_i,\\
  &\theta=\beta_i,\quad r>d_i:
    &&u_\theta(r,\beta_i)=0,\\
  &\theta=\beta_j,\quad r>0,\quad j\ne i:
    &&u_\theta(r,\beta_j)=0.
\end{aligned}
\end{equation}
Отже, умова симетрії для $u_\theta$ замінюється умовою сталого відривного
напруження лише на ділянці $0<r<d_i$ активного напрямку $\theta=\beta_i$.

На проміжних відстанях $d_i\ll r\ll l$ розв'язок повинен переходити в
асимптотику задачі про міжфазну тріщину без зони передруйнування~\cite{NazarenkoKipnis2022Stress}. У
внутрішньому масштабі ця умова записується як
\begin{equation}
\label{eq:matched-far-field}
  \sigma_\theta(r,\beta_i)
    =CQ_i r^\lambda+o\!\left(r^{-1}\right),
    \qquad r\to\infty.
\end{equation}
Тут $Q_i$ --- функція кута і пружних характеристик матеріалів, а
$\lambda\in(-1,0)$ --- показник сингулярності кутової точки після
утворення міжфазної тріщини. Під $\lambda$ розуміємо найменший
дійсний корінь рівняння $D(-1-\lambda)=0$ у $(-1,0)$, неперервно
пов'язаний із фізичною гілкою. Формули для $D$ і $Q_i$ наведено в
додатку~\ref{app:asymptotics}.

Для подальшого застосування методу Вінера--Гопфа використовуємо принцип
суперпозиції так само, як у спорідненій задачі~\cite{KaminskyDudykPolishchuk2024}.
Із повного поля віднімаємо відоме поле задачі без зони передруйнування та
вводимо коригувальне нормальне напруження
\begin{equation*}
  \widehat{\sigma}_i(r)
  =\sigma_\theta(r,\beta_i)-CQ_i r^\lambda .
\end{equation*}
Тоді з умови зрощування випливає
\begin{equation*}
  \widehat{\sigma}_i(r)=o(r^{-1}),
  \qquad r\to\infty,
\end{equation*}
а на ділянці зони $0<r<d_i$
\begin{equation*}
  \widehat{\sigma}_i(r)
  =\sigma_i-CQ_i r^\lambda.
\end{equation*}
Поле без зони задовольняє умову симетрії $u_\theta=0$ на бісектрисі,
тому віднімання цього поля не змінює трансформанту градієнта
$u_\theta$, що використовується нижче. Отже, подальша функціональна
задача формулюється для коригувального поля; відомий степеневий доданок
входить до її заданої частини.

Зона передруйнування може виникнути в матеріалі~$i$ лише за умови
розтягувального нормального напруження на відповідній бісектрисі. З
формули~\eqref{eq:matched-far-field} випливає критерій фізичної допустимості
\begin{equation*}
  CQ_i>0.
\end{equation*}
Отже, модель потребує одночасного виконання двох нерівностей:
$Cg_2<0$ для контакту берегів міжфазної тріщини та $CQ_i>0$ для
утворення зони в матеріалі~$i$. Для досліджених нижче наборів параметрів
числовий аналіз показує, що ці умови можуть одночасно виконуватися лише
для одного з двох матеріалів; саме в ньому фізично допустиме утворення
зони передруйнування.

\subsection{Асимптотика біля кінця зони}
\label{subsec:zone-tip-asymptotics}

Біля зовнішнього кінця лінії розриву напруження і градієнт нормального
переміщення мають стандартну кореневу асимптотику однорідного матеріалу:
\begin{equation}
\label{eq:zone-tip-asymptotics}
\begin{aligned}
  &\theta=\beta_i,\quad r\to d_i+0:
  &&\sigma_\theta(r,\beta_i)
    \sim \frac{k_i}{\sqrt{2\pi(r-d_i)}},\\[3pt]
  &\theta=\beta_i,\quad r\to d_i-0:
  &&\frac{\partial u_\theta}{\partial r}(r,\beta_i)
    \sim -\frac{2(1-\nu_i^2)}{E_i}
    \frac{k_i}{\sqrt{2\pi(d_i-r)}}.
\end{aligned}
\end{equation}
Коефіцієнт $k_i$ є локальним коефіцієнтом інтенсивності напружень у
кінці зони. Для замикання крайової задачі використовуємо стандартну для
цієї моделі умову усунення кореневого сингулярного доданка
\cite{KaminskyDudykPolishchuk2024}
\begin{equation*}
  k_i=0.
\end{equation*}
Ця умова зануляє коефіцієнт обернено-кореневого сингулярного члена і
далі використовується для визначення рівноважної довжини $d_i$; вона не
вимагає, щоб обмежена частина напруження в кінці зони дорівнювала нулю.